\documentclass[10pt,a4paper]{article}
\usepackage{xeCJK}
\usepackage{fontspec}
\usepackage{booktabs}
\usepackage{array}
\usepackage{xcolor}
\usepackage{geometry}
\usepackage[protrusion=true]{microtype}
\geometry{margin=1.7cm}
\linespread{1.05}

\setmainfont{Baskerville}
\setCJKmainfont{Songti SC}
\newfontfamily{\starfont}{Songti SC}

\definecolor{wrongred}{RGB}{198,40,40}
\definecolor{rulegray}{RGB}{190,190,190}

\newcommand{\checkmarkbox}{\fbox{\rule{0pt}{0.75ex}\rule{0.75ex}{0pt}}}
\newcommand{\STAR}{\textcolor{wrongred}{\starfont ★}}
% \fa{题号}{答案}{解析}
\newcommand{\fa}[3]{\noindent\textbf{#1.}~\textbf{#2}\par\nobreak\vspace{0.05em}\hspace*{1.5em}{\small #3}\par\vspace{0.42em}}

\begin{document}
\pagestyle{plain}

\begin{center}
{\LARGE\bfseries 时态语态练习1 \quad 参考答案与解析}\\[0.25em]
{\large 智学网 AI 错题本 · 原错题 1--18（18 题、21 空；按题计分）}\\[0.22em]
{\small 2029届高一英语 · 解析以"时间状语定调 + 主谓一致"为主线}
\end{center}

\vspace{0.3em}
\noindent{\small 说明：题号沿用原错题本。第 9 题答案册作 \emph{were working}，本表按标准语法取 \emph{was working}。}

\vspace{0.45em}
\fa{1}{will be sitting}{this time next month（下个月的这个时候）$\to$ 将来进行时。}
\fa{2}{had been}{"It was the second time (that) sb had done"：主句过去，从句用过去完成。}
\fa{3}{has made}{按本练习答案口径，用现在完成强调“已下定决心”；一般条件表达中 makes 也自然，题干未完全排除。}
\fa{4}{were writing; had left}{when 提供背景 $\to$ 主句用过去进行 were writing；书在"去取"之前就被落在办公室 $\to$ 过去完成 had left。}
\fa{5}{will have worked}{到年底将累计任教 30 年，预期用将来完成时；若强调届时仍在持续，可用 will have been working。}
\fa{6}{is always making}{Betty 作主语；always + 进行时可表“老是……”的抱怨语气。}
\fa{7}{have seen}{Recent years 作复数主语，指近些年；用现在完成时 have seen（years 见证了上海的巨大变化）。}
\fa{8}{has been tried}{since then $\to$ 现在完成；means 作主语是单数，且与 try 是被动关系。}
\fa{9}{was working}{when 从句背景 $\to$ 过去进行；everybody 作主语是单数，用 was。}
\fa{10}{was staring}{"我进来时他正盯着空中发呆" $\to$ 过去进行时。}
\fa{11}{have been painting}{all day（一整天，持续到现在）$\to$ 现在完成进行时。}
\fa{12}{was writing}{last year 那段时间"正在写" $\to$ 过去进行时。}
\fa{13}{didn't recognize; have just had; am wearing}{"认出来"发生在说话瞬间 $\to$ 一般过去；"刚刚理过发"$\to$ 现在完成；"戴着新眼镜"（当前状态）$\to$ 现在进行。}
\fa{14}{have published}{in the last three years $\to$ 现在完成时。}
\fa{15}{worked}{in the 1990s 指已结束的九十年代，故用一般过去时；年代复数不加撇号。}
\fa{16}{am helping}{these days 表临时性 $\to$ 现在进行时。}
\fa{17}{was leaving / would leave}{主句 He said 是过去，the next day 表"过去的将来" $\to$ 过去将来时。}
\fa{18}{was thinking}{按原练习口径，was thinking 可作委婉、试探的过去表达，不必然表示计划落空；若从现在出发考虑，也可说 am thinking，题干缺少明确过去时间锚点。}

\vspace{0.3em}
\noindent{\color{rulegray}\rule{\textwidth}{0.5pt}}
\vspace{0.3em}
\noindent\textbf{复核记录}\hfill\small 完成重做后填写

\vspace{0.3em}
\noindent\begin{tabular}{@{}p{0.31\textwidth}@{}p{0.34\textwidth}@{}p{0.35\textwidth}@{}}
重做得分：\underline{\hspace{1.2cm}}\,/\,18
& 仍错误的题号：\underline{\hspace{3.2cm}}
& 下次复习日期：\underline{\hspace{2.0cm}} \\
\end{tabular}

\vspace{0.3em}
\noindent 本轮易错点：\checkmarkbox\ 完成时/进行时\quad\checkmarkbox\ 时间状语定调\quad\checkmarkbox\ 被动与主谓一致\quad\checkmarkbox\ 过去将来

\vspace{0.3em}
\noindent 复习状态：\checkmarkbox\ 已掌握\quad\checkmarkbox\ 需要巩固\quad\checkmarkbox\ 需重点复习

\end{document}
